% Source: https://artofproblemsolving.com/community/c3h1607743_famous_mathematicians
% Stable identity: aops-topic-1607743-q01
There are ten different letters $E,U,C,L,I,D,G,A,S,P$, so every digit $0$--$9$ is used exactly once. Write $c_1,\dots,c_5$ for the carries out of the units, tens, hundreds, thousands and ten-thousands columns.

\textbf{The column equations.} Reading the columns from the right,
\[
\begin{aligned}
&\text{units: } D+S=L+10c_1, &&\text{tens: } I+S+c_1=A+10c_2,\\
&\text{hundreds: } L+U+c_2=C+10c_3, &&\text{thousands: } C+A+c_3=S+10c_4,\\
&\text{ten-thousands: } U+G+c_4=A+10c_5, &&\text{hundred-thousands: } E+c_5=P .
\end{aligned}
\]

\textbf{First deductions.}
\begin{itemize}
\item Since $P\ne E$, the last equation forces $c_5=1$ and $P=E+1$.
\item Adding the tens and thousands equations and cancelling $A$ and $S$ gives $I+C+c_1+c_3=10(c_2+c_4)$. The left side is positive (as $I\ne C$) and at most $9+8+1+1=19$, so $c_2+c_4=1$ and
\[
I+C+c_1+c_3=10 .
\]
\item The ten-thousands equation becomes $U+G+c_4=A+10$, so $U+G\ge 9+A$.
\end{itemize}
These equations cut the possibilities down sharply. Running through the remaining cases (the carries $c_1,c_3$ and the choice of which of $c_2$, $c_4$ is $1$, then the digits allowed by each equation, always keeping the ten letters distinct) leaves exactly one assignment:
\[
E=2,\ U=8,\ C=9,\ L=1,\ I=0,\ D=7,\ G=6,\ A=5,\ S=4,\ P=3 ,
\]
with carries $c_1=1$, $c_2=0$, $c_3=0$, $c_4=1$, $c_5=1$ (indeed $I+C+c_1+c_3=0+9+1+0=10$). that is,
\[
\boxed{289107+65844=354951}.
\]
Check by columns: $7+4=11$ (write $1=L$, carry $1$); $0+4+1=5=A$; $1+8=9=C$; $9+5=14$ (write $4=S$, carry $1$); $8+6+1=15$ (write $5=A$, carry $1$); $2+1=3=P$. A complete search over all assignments of distinct digits confirms that this solution is unique.
